\section{Results}
\subsection{Accuracy Comparison}
Table~\ref{tab:full_results} presents the full results for all ranks. Each entry is the mean of three runs (standard deviation $<0.2$ for all except CoLA where it is $<0.5$).

\begin{table*}[t]
\centering
\caption{GLUE validation scores for SVD-Lin with different ranks. Full BERT baseline is rank $\infty$.}
\label{tab:full_results}
\begin{tabular}{lcccc}
\toprule
\textbf{Task} & \textbf{Rank 64} & \textbf{Rank 128} & \textbf{Rank 256} & \textbf{Full BERT} \\
\midrule
CoLA (MCC)       & 58.5 & 59.8 & 60.2 & 60.9 \\
SST-2 (Acc)      & 91.4 & 92.2 & 92.7 & 93.0 \\
MRPC (Acc)       & 87.3 & 88.1 & 88.7 & 89.2 \\
QQP (Acc)        & 90.7 & 91.2 & 91.5 & 91.8 \\
QNLI (Acc)       & 89.2 & 90.3 & 90.9 & 91.4 \\
RTE (Acc)        & 67.5 & 68.3 & 69.0 & 69.7 \\
MNLI-m (Acc)     & 82.3 & 83.1 & 83.8 & 84.3 \\
STS-B (Pearson)  & 87.6 & 88.4 & 88.8 & 89.0 \\
\bottomrule
\end{tabular}
\end{table*}

\textbf{Observations:}
- Rank 64 achieves at least 96\% of full accuracy on all tasks (e.g., SST-2: 91.4 vs 93.0 = 98.3\%).
- Rank 128 narrows the gap to within 0.5\% absolute (e.g., MNLI: 83.1 vs 84.3 = 98.6\%).
- Rank 256 is nearly indistinguishable from full BERT (maximum difference 0.3\%).
- CoLA is the most sensitive to compression, likely because grammaticality judgments require fine-grained token interactions.

\subsection{Efficiency Gains}
Table~\ref{tab:efficiency} shows the relative training time and memory usage normalized to full BERT. Measurements are for SST-2 with batch size 16, sequence length 128.

\begin{table}[h]
\centering
\caption{Relative efficiency on M1 Mac (SST-2).}
\label{tab:efficiency}
\begin{tabular}{lccc}
\toprule
\textbf{Rank} & \textbf{Time/epoch (s)} & \textbf{Peak Memory (MB)} & \textbf{Relative Speed} \\
\midrule
64   & 243 & 1,120 & 2.31$\times$ \\
128  & 352 & 1,720 & 1.60$\times$ \\
256  & 458 & 2,450 & 1.23$\times$ \\
Full & 562 & 3,140 & 1.00$\times$ \\
\bottomrule
\end{tabular}
\end{table}

Rank 64 reduces memory by nearly 3× and time by 2.3×. Even rank 128 provides substantial savings (1.6× speedup). The memory reduction is especially important for deployment on edge devices.

\subsection{Effect of Sequence Length}
We artificially extended sequences by repeating tokens (with attention mask) to lengths 256, 512, and 1024, using a fixed rank $r=128$. Figure~\ref{fig:scaling} shows the scaling of forward pass time. Full attention grows quadratically, while SVD-Lin grows linearly.

% \begin{figure}[h]
% \centering
% % \includegraphics[width=0.48\textwidth]{scaling.png}
% \caption{Forward pass time vs. sequence length (log-log). SVD-Lin (r=128) scales as $O(N)$, full attention as $O(N^2)$.}
% \label{fig:scaling}
% \end{figure}

\subsection{Ablation: SVD vs. Random Projection}
To isolate the benefit of data-dependent SVD, we compared against a variant where $U_r$ is replaced by a random orthonormal matrix (fixed per layer). Table~\ref{tab:ablation} shows results on SST-2 and MNLI for rank 128.

\begin{table}[h]
\centering
\caption{SVD vs. random projection (rank 128).}
\label{tab:ablation}
\begin{tabular}{lcc}
\toprule
\textbf{Method} & \textbf{SST-2 Acc} & \textbf{MNLI-m Acc} \\
\midrule
SVD-Lin (ours) & 92.2 & 83.1 \\
Random projection & 91.5 & 82.3 \\
Full BERT & 93.0 & 84.3 \\
\bottomrule
\end{tabular}
\end{table}

SVD-Lin outperforms the random projection by 0.7\% on SST-2 and 0.8\% on MNLI, confirming that adaptivity matters.

\subsection{Sensitivity to Rank}
We swept ranks from 16 to 512 on SST-2 (Figure~\ref{fig:sensitivity}). Accuracy improves rapidly up to rank 64, then saturates. This suggests that the effective rank of BERT's representations is around 64–128, aligning with prior observations on language model intrinsic dimension.

% \begin{figure}[h]
% \centering
% \includegraphics[width=0.48\textwidth]{sensitivity.png}
% \caption{Accuracy on SST-2 as a function of rank.}
% \label{fig:sensitivity}
% \end{figure}